% Source: backpropagation_example.tex
\begin{notechapter}
\chapter{Backpropagation: A Worked Example}\label{ch:backprop-worked-example}

\section{Network Architecture}

Consider a neural network with the following architecture:
\begin{itemize}
    \item Input layer: \( \mathbf{x} \in \mathbb{R}^{2 \times 1} \),
    \item Hidden layer: \( \mathbf{h} \in \mathbb{R}^{3 \times 1} \),
    \item Output layer: \( \mathbf{y} \in \mathbb{R}^{2 \times 1} \).
\end{itemize}

The network has the following parameters:
\begin{itemize}
    \item Weight matrix \( W_1 \in \mathbb{R}^{3 \times 2} \) for the 
          input-to-hidden layer,
    \item Bias vector \( \mathbf{b}_1 \in \mathbb{R}^{3 \times 1} \) 
          for the hidden layer,
    \item Weight matrix \( W_2 \in \mathbb{R}^{2 \times 3} \) for the 
          hidden-to-output layer,
    \item Bias vector \( \mathbf{b}_2 \in \mathbb{R}^{2 \times 1} \) 
          for the output layer.
\end{itemize}

The activation function for the hidden layer is \( \sigma \), 
applied element-wise.

\section{Forward Pass}

\subsection{Hidden Layer}

The hidden layer activations \( \mathbf{h} \) are computed as:
\[
\mathbf{h} = \sigma(W_1 \mathbf{x} + \mathbf{b}_1),
\]
where:
\begin{itemize}
    \item \( W_1 \mathbf{x} + \mathbf{b}_1 \) is the pre-activation 
          vector of size \( 3 \times 1 \),
    \item \( \sigma \) is the activation function (e.g., ReLU, sigmoid), 
          applied element-wise.
\end{itemize}

\subsubsection{Explicit Matrices}

Let:
\[
W_1 = \begin{bmatrix}
w_{11} & w_{12} \\
w_{21} & w_{22} \\
w_{31} & w_{32}
\end{bmatrix}, \quad
\mathbf{x} = \begin{bmatrix}
x_1 \\
x_2
\end{bmatrix}, \quad
\mathbf{b}_1 = \begin{bmatrix}
b_{11} \\
b_{21} \\
b_{31}
\end{bmatrix}.
\]

The pre-activation vector \( \mathbf{z} = W_1 \mathbf{x} + \mathbf{b}_1 \) is:
\[
\mathbf{z} = \begin{bmatrix}
z_1 \\
z_2 \\
z_3
\end{bmatrix} = \begin{bmatrix}
w_{11}x_1 + w_{12}x_2 + b_{11} \\
w_{21}x_1 + w_{22}x_2 + b_{21} \\
w_{31}x_1 + w_{32}x_2 + b_{31}
\end{bmatrix}.
\]

After applying the activation function \( \sigma \):
\[
\mathbf{h} = \begin{bmatrix}
\sigma(z_1) \\
\sigma(z_2) \\
\sigma(z_3)
\end{bmatrix}.
\]

\subsection{Output Layer}

The output \( \mathbf{y} \) is computed as:
\[
\mathbf{y} = W_2 \mathbf{h} + \mathbf{b}_2,
\]
where:
\begin{itemize}
    \item \( W_2 \mathbf{h} + \mathbf{b}_2 \) is the output vector 
          of size \( 2 \times 1 \).
\end{itemize}

\subsubsection{Explicit Matrices}

Let:
\[
W_2 = \begin{bmatrix}
w_{11}^{\prime} & w_{12}^{\prime} & w_{13}^{\prime} \\
w_{21}^{\prime} & w_{22}^{\prime} & w_{23}^{\prime}
\end{bmatrix}, \quad
\mathbf{h} = \begin{bmatrix}
h_1 \\
h_2 \\
h_3
\end{bmatrix}, \quad
\mathbf{b}_2 = \begin{bmatrix}
b_{11}^{\prime} \\
b_{21}^{\prime}
\end{bmatrix}.
\]

Then:
\[
W_2 \mathbf{h} + \mathbf{b}_2 = \begin{bmatrix}
w_{11}^{\prime}h_1 + w_{12}^{\prime}h_2 + w_{13}^{\prime}h_3 + b_{11}^{\prime} \\
w_{21}^{\prime}h_1 + w_{22}^{\prime}h_2 + w_{23}^{\prime}h_3 + b_{21}^{\prime}
\end{bmatrix}.
\]

Thus:
\[
\mathbf{y} = \begin{bmatrix}
y_1 \\
y_2
\end{bmatrix} = \begin{bmatrix}
w_{11}^{\prime}h_1 + w_{12}^{\prime}h_2 + w_{13}^{\prime}h_3 + b_{11}^{\prime} \\
w_{21}^{\prime}h_1 + w_{22}^{\prime}h_2 + w_{23}^{\prime}h_3 + b_{21}^{\prime}
\end{bmatrix}.
\]

\section{Loss Function}

Let \( \mathcal{L} \) be the loss function, which is a scalar 
function of the output \( \mathbf{y} \). For example, for mean 
squared error (MSE) loss:
\[
\mathcal{L} = \frac{1}{2} \|\mathbf{y} - \mathbf{y}_{\text{target}}\|^2,
\]
where \( \mathbf{y}_{\text{target}} \) is the target output.

\section{Backward Pass}

The backward pass computes the gradients of the loss \( \mathcal{L} \) 
with respect to the network parameters \( W_1, \mathbf{b}_1, W_2, \mathbf{b}_2 \) 
using the chain rule.

\subsection{Gradient with Respect to \( W_2 \)}

The gradient of \( \mathcal{L} \) with respect to \( W_2 \) is:
\[
\frac{\partial \mathcal{L}}{\partial W_2} = \frac{\partial \mathcal{L}}{\partial \mathbf{y}} 
\cdot \frac{\partial \mathbf{y}}{\partial W_2}.
\]

Here:
\begin{itemize}
    \item \( \frac{\partial \mathcal{L}}{\partial \mathbf{y}} \) is a \( 1 \times 2 \) 
          row vector (gradient of the scalar loss with respect to the output \( \mathbf{y} \)),
    \item \( \frac{\partial \mathbf{y}}{\partial W_2} \) is a \( 2 \times 2 \times 3 \) 
          tensor (derivative of the vector \( \mathbf{y} \) with respect to the matrix \( W_2 \)).
\end{itemize}

\subsubsection{Explicit Tensor \( \frac{\partial \mathbf{y}}{\partial W_2} \)}

The tensor \( \frac{\partial \mathbf{y}}{\partial W_2} \) has dimensions \( 2 \times 2 \times 3 \), 
where:
\begin{itemize}
    \item The first dimension corresponds to the elements of \( \mathbf{y} \) (\( y_1 \) and \( y_2 \)),
    \item The second dimension corresponds to the rows of \( W_2 \) (\( 2 \) rows),
    \item The third dimension corresponds to the columns of \( W_2 \) (\( 3 \) columns).
\end{itemize}

For example:
\[
\frac{\partial y_1}{\partial W_2} = \begin{bmatrix}
\frac{\partial y_1}{\partial w_{11}^{\prime}} & \frac{\partial y_1}{\partial w_{12}^{\prime}} & \frac{\partial y_1}{\partial w_{13}^{\prime}} \\
\frac{\partial y_1}{\partial w_{21}^{\prime}} & \frac{\partial y_1}{\partial w_{22}^{\prime}} & \frac{\partial y_1}{\partial w_{23}^{\prime}}
\end{bmatrix} = \begin{bmatrix}
h_1 & h_2 & h_3 \\
0 & 0 & 0
\end{bmatrix},
\]
and
\[
\frac{\partial y_2}{\partial W_2} = \begin{bmatrix}
0 & 0 & 0 \\
h_1 & h_2 & h_3
\end{bmatrix}.
\]

Thus, the full tensor \( \frac{\partial \mathbf{y}}{\partial W_2} \) is:
\[
\frac{\partial \mathbf{y}}{\partial W_2} = \begin{bmatrix}
\begin{bmatrix}
h_1 & h_2 & h_3 \\
0 & 0 & 0
\end{bmatrix},
\begin{bmatrix}
0 & 0 & 0 \\
h_1 & h_2 & h_3
\end{bmatrix}
\end{bmatrix}.
\]

\subsection{Gradient with Respect to \( \mathbf{b}_2 \)}

The gradient of \( \mathcal{L} \) with respect to \( \mathbf{b}_2 \) is:
\[
\frac{\partial \mathcal{L}}{\partial \mathbf{b}_2} = \frac{\partial \mathcal{L}}{\partial \mathbf{y}}.
\]
This is because \( \mathbf{y} = W_2 \mathbf{h} + \mathbf{b}_2 \), so \( \frac{\partial \mathbf{y}}{\partial \mathbf{b}_2} = I \), 
where \( I \) is the identity matrix.

\subsection{Gradient with Respect to \( \mathbf{h} \)}

The gradient of \( \mathcal{L} \) with respect to \( \mathbf{h} \) is:
\[
\frac{\partial \mathcal{L}}{\partial \mathbf{h}} = \frac{\partial \mathcal{L}}{\partial \mathbf{y}} \cdot W_2.
\]
Here:
\begin{itemize}
    \item \( \frac{\partial \mathcal{L}}{\partial \mathbf{y}} \) is a \( 1 \times 2 \) row vector,
    \item \( W_2 \) is a \( 2 \times 3 \) matrix,
    \item The product \( \frac{\partial \mathcal{L}}{\partial \mathbf{y}} \cdot W_2 \) is a \( 1 \times 3 \) 
          row vector, matching the dimensions of \( \mathbf{h} \).
\end{itemize}

\subsection{Gradient with Respect to \( W_1 \)}

The gradient of \( \mathcal{L} \) with respect to \( W_1 \) is:
\[
\frac{\partial \mathcal{L}}{\partial W_1} = \frac{\partial \mathcal{L}}{\partial \mathbf{h}} 
\cdot \frac{\partial \mathbf{h}}{\partial W_1}.
\]

Here:
\begin{itemize}
    \item \( \frac{\partial \mathcal{L}}{\partial \mathbf{h}} \) is a \( 1 \times 3 \) row vector,
    \item \( \frac{\partial \mathbf{h}}{\partial W_1} \) is a \( 3 \times 3 \times 2 \) tensor 
          (derivative of the vector \( \mathbf{h} \) with respect to the matrix \( W_1 \)).
\end{itemize}

\subsubsection{Explicit Tensor \( \frac{\partial \mathbf{h}}{\partial W_1} \)}

The tensor \( \frac{\partial \mathbf{h}}{\partial W_1} \) has dimensions \( 3 \times 3 \times 2 \), 
where:
\begin{itemize}
    \item The first dimension corresponds to the elements of \( \mathbf{h} \) (\( h_1, h_2, h_3 \)),
    \item The second dimension corresponds to the rows of \( W_1 \) (\( 3 \) rows),
    \item The third dimension corresponds to the columns of \( W_1 \) (\( 2 \) columns).
\end{itemize}

For example:
\[
\frac{\partial h_1}{\partial W_1} = \begin{bmatrix}
\frac{\partial h_1}{\partial w_{11}} & \frac{\partial h_1}{\partial w_{12}} \\
\frac{\partial h_1}{\partial w_{21}} & \frac{\partial h_1}{\partial w_{22}} \\
\frac{\partial h_1}{\partial w_{31}} & \frac{\partial h_1}{\partial w_{32}}
\end{bmatrix} = \begin{bmatrix}
\sigma^{\prime}(z_1) x_1 & \sigma^{\prime}(z_1) x_2 \\
0 & 0 \\
0 & 0
\end{bmatrix},
\]
where \( z_1 = w_{11}x_1 + w_{12}x_2 + b_{11} \).

Similarly:
\[
\frac{\partial h_2}{\partial W_1} = \begin{bmatrix}
0 & 0 \\
\sigma^{\prime}(z_2) x_1 & \sigma^{\prime}(z_2) x_2 \\
0 & 0
\end{bmatrix}, \quad
\frac{\partial h_3}{\partial W_1} = \begin{bmatrix}
0 & 0 \\
0 & 0 \\
\sigma^{\prime}(z_3) x_1 & \sigma^{\prime}(z_3) x_2
\end{bmatrix}.
\]

Thus, the full tensor \( \frac{\partial \mathbf{h}}{\partial W_1} \) is:
\[
\frac{\partial \mathbf{h}}{\partial W_1} = \begin{bmatrix}
\begin{bmatrix}
\sigma^{\prime}(z_1) x_1 & \sigma^{\prime}(z_1) x_2 \\
0 & 0 \\
0 & 0
\end{bmatrix},
\begin{bmatrix}
0 & 0 \\
\sigma^{\prime}(z_2) x_1 & \sigma^{\prime}(z_2) x_2 \\
0 & 0
\end{bmatrix},
\begin{bmatrix}
0 & 0 \\
0 & 0 \\
\sigma^{\prime}(z_3) x_1 & \sigma^{\prime}(z_3) x_2
\end{bmatrix}
\end{bmatrix}.
\]

\subsection{Gradient with Respect to \( \mathbf{b}_1 \)}

The gradient of \( \mathcal{L} \) with respect to \( \mathbf{b}_1 \) is:
\[
\frac{\partial \mathcal{L}}{\partial \mathbf{b}_1} = \frac{\partial \mathcal{L}}{\partial \mathbf{h}} 
\odot \sigma^{\prime}(\mathbf{z}),
\]
where:
\begin{itemize}
    \item \( \frac{\partial \mathcal{L}}{\partial \mathbf{h}} \) is a \( 1 \times 3 \) row vector,
    \item \( \sigma^{\prime}(\mathbf{z}) \) is a \( 3 \times 1 \) column vector,
    \item The Hadamard product \( \frac{\partial \mathcal{L}}{\partial \mathbf{h}} \odot \sigma^{\prime}(\mathbf{z}) \) 
          is a \( 3 \times 1 \) column vector, matching the dimensions of \( \mathbf{b}_1 \).
\end{itemize}

\subsubsection{Derivation of \( \frac{\partial \mathbf{h}}{\partial \mathbf{b}_1} \)}

The gradient \( \frac{\partial \mathbf{h}}{\partial \mathbf{b}_1} \) is derived using the chain rule:
\[
\frac{\partial \mathbf{h}}{\partial \mathbf{b}_1} = \frac{\partial \mathbf{h}}{\partial \mathbf{z}} \cdot \frac{\partial \mathbf{z}}{\partial \mathbf{b}_1}.
\]

Here:
\begin{itemize}
    \item \( \frac{\partial \mathbf{h}}{\partial \mathbf{z}} \) is the derivative of the activation function \( \sigma \) applied element-wise to \( \mathbf{z} \). Since \( \sigma \) is applied element-wise, \( \frac{\partial \mathbf{h}}{\partial \mathbf{z}} \) is a diagonal matrix where each diagonal entry is \( \sigma^{\prime}(z_i) \), the derivative of \( \sigma \) at \( z_i \). 

    To derive this, note that:
    \[
    \mathbf{h} = \begin{bmatrix}
    \sigma(z_1) \\
    \sigma(z_2) \\
    \sigma(z_3)
    \end{bmatrix},
    \]
    and the derivative of \( \mathbf{h} \) with respect to \( \mathbf{z} \) is:
    \[
    \frac{\partial \mathbf{h}}{\partial \mathbf{z}} = \begin{bmatrix}
    \frac{\partial \sigma(z_1)}{\partial z_1} & \frac{\partial \sigma(z_1)}{\partial z_2} & \frac{\partial \sigma(z_1)}{\partial z_3} \\
    \frac{\partial \sigma(z_2)}{\partial z_1} & \frac{\partial \sigma(z_2)}{\partial z_2} & \frac{\partial \sigma(z_2)}{\partial z_3} \\
    \frac{\partial \sigma(z_3)}{\partial z_1} & \frac{\partial \sigma(z_3)}{\partial z_2} & \frac{\partial \sigma(z_3)}{\partial z_3}
    \end{bmatrix}.
    \]

    Since \( \sigma \) is applied element-wise, the off-diagonal terms are zero, and the diagonal terms are \( \sigma^{\prime}(z_i) \). Thus:
    \[
    \frac{\partial \mathbf{h}}{\partial \mathbf{z}} = \begin{bmatrix}
    \sigma^{\prime}(z_1) & 0 & 0 \\
    0 & \sigma^{\prime}(z_2) & 0 \\
    0 & 0 & \sigma^{\prime}(z_3)
    \end{bmatrix}.
    \]

    \item \( \frac{\partial \mathbf{z}}{\partial \mathbf{b}_1} \) is the derivative of \( \mathbf{z} \) with respect to \( \mathbf{b}_1 \). Since \( \mathbf{z} = W_1 \mathbf{x} + \mathbf{b}_1 \), the derivative of \( \mathbf{z} \) with respect to \( \mathbf{b}_1 \) is the identity matrix \( I \) of size \( 3 \times 3 \). To derive this, note that:
    \[
    \mathbf{z} = \begin{bmatrix}
    z_1 \\
    z_2 \\
    z_3
    \end{bmatrix} = \begin{bmatrix}
    w_{11}x_1 + w_{12}x_2 + b_{11} \\
    w_{21}x_1 + w_{22}x_2 + b_{21} \\
    w_{31}x_1 + w_{32}x_2 + b_{31}
    \end{bmatrix}.
    \]

    The derivative of \( \mathbf{z} \) with respect to \( \mathbf{b}_1 \) is:
    \[
    \frac{\partial \mathbf{z}}{\partial \mathbf{b}_1} = \begin{bmatrix}
    \frac{\partial z_1}{\partial b_{11}} & \frac{\partial z_1}{\partial b_{21}} & \frac{\partial z_1}{\partial b_{31}} \\
    \frac{\partial z_2}{\partial b_{11}} & \frac{\partial z_2}{\partial b_{21}} & \frac{\partial z_2}{\partial b_{31}} \\
    \frac{\partial z_3}{\partial b_{11}} & \frac{\partial z_3}{\partial b_{21}} & \frac{\partial z_3}{\partial b_{31}}
    \end{bmatrix}.
    \]

    Since \( z_i \) depends only on \( b_{i1} \) (e.g., \( z_1 = w_{11}x_1 + w_{12}x_2 + b_{11} \)), the off-diagonal terms are zero, and the diagonal terms are 1. Thus:
    \[
    \frac{\partial \mathbf{z}}{\partial \mathbf{b}_1} = \begin{bmatrix}
    1 & 0 & 0 \\
    0 & 1 & 0 \\
    0 & 0 & 1
    \end{bmatrix} = I.
    \]
\end{itemize}

Combining these results, we have:
\[
\frac{\partial \mathbf{h}}{\partial \mathbf{b}_1} = \frac{\partial \mathbf{h}}{\partial \mathbf{z}} \cdot \frac{\partial \mathbf{z}}{\partial \mathbf{b}_1} = \begin{bmatrix}
\sigma^{\prime}(z_1) & 0 & 0 \\
0 & \sigma^{\prime}(z_2) & 0 \\
0 & 0 & \sigma^{\prime}(z_3)
\end{bmatrix} \cdot \begin{bmatrix}
1 & 0 & 0 \\
0 & 1 & 0 \\
0 & 0 & 1
\end{bmatrix} = 
\begin{bmatrix}
\sigma^{\prime}(z_1) & 0 & 0 \\
0 & \sigma^{\prime}(z_2) & 0 \\
0 & 0 & \sigma^{\prime}(z_3)
\end{bmatrix}.
\]

Thus, \( \frac{\partial \mathbf{h}}{\partial \mathbf{b}_1} \) is a \( 3 \times 3 \) diagonal matrix where each diagonal entry is the derivative of the activation function \( \sigma \) evaluated at the corresponding element of \( \mathbf{z} \).

\section{Summary of Gradients}

\subsection{Output Layer Gradients}

The gradients of the loss \( \mathcal{L} \) with respect to the 
output layer parameters \( W_2 \) and \( \mathbf{b}_2 \) are:

\begin{itemize}
    \item Gradient with respect to \( W_2 \):
    \[
    \frac{\partial \mathcal{L}}{\partial W_2} = 
    \left( \frac{\partial \mathcal{L}}{\partial \mathbf{y}} \right)^T 
    \cdot \mathbf{h}^T.
    \]
    This is derived as follows:
    \begin{itemize}
        \item The chain rule gives:
        \[
        \frac{\partial \mathcal{L}}{\partial W_2} = 
        \frac{\partial \mathcal{L}}{\partial \mathbf{y}} 
        \cdot \frac{\partial \mathbf{y}}{\partial W_2}.
        \]
        \item Since \( \mathbf{y} = W_2 \mathbf{h} + \mathbf{b}_2 \), the 
        derivative \( \frac{\partial \mathbf{y}}{\partial W_2} \) is a 
        \( 2 \times 2 \times 3 \) tensor. To compute it, we first write:
        \[
        \frac{\partial \mathcal{L}}{\partial W_2} = 
        \frac{\partial \mathcal{L}}{\partial y_1} \cdot 
        \frac{\partial y_1}{\partial W_2} + 
        \frac{\partial \mathcal{L}}{\partial y_2} \cdot 
        \frac{\partial y_2}{\partial W_2}.
        \]
        \item The partial derivatives \( \frac{\partial y_1}{\partial W_2} \) 
        and \( \frac{\partial y_2}{\partial W_2} \) are matrices of size 
        \( 2 \times 3 \), where:
        \[
        \frac{\partial y_1}{\partial W_2} = 
        \begin{bmatrix}
        h_1 & h_2 & h_3 \\
        0 & 0 & 0
        \end{bmatrix}, \quad
        \frac{\partial y_2}{\partial W_2} = 
        \begin{bmatrix}
        0 & 0 & 0 \\
        h_1 & h_2 & h_3
        \end{bmatrix}.
        \]
        \item Substituting these into the expression for 
        \( \frac{\partial \mathcal{L}}{\partial W_2} \), we get:
        \[
        \frac{\partial \mathcal{L}}{\partial W_2} = 
        \frac{\partial \mathcal{L}}{\partial y_1} \cdot 
        \begin{bmatrix}
        h_1 & h_2 & h_3 \\
        0 & 0 & 0
        \end{bmatrix} + 
        \frac{\partial \mathcal{L}}{\partial y_2} \cdot 
        \begin{bmatrix}
        0 & 0 & 0 \\
        h_1 & h_2 & h_3
        \end{bmatrix}.
        \]
        \item This results in:
        \[
        \frac{\partial \mathcal{L}}{\partial W_2} = 
        \begin{bmatrix}
        \frac{\partial \mathcal{L}}{\partial y_1} h_1 & 
        \frac{\partial \mathcal{L}}{\partial y_1} h_2 & 
        \frac{\partial \mathcal{L}}{\partial y_1} h_3 \\
        \frac{\partial \mathcal{L}}{\partial y_2} h_1 & 
        \frac{\partial \mathcal{L}}{\partial y_2} h_2 & 
        \frac{\partial \mathcal{L}}{\partial y_2} h_3
        \end{bmatrix}.
        \]
        \item This can be written compactly as:
        \[
        \frac{\partial \mathcal{L}}{\partial W_2} = 
        \left( \frac{\partial \mathcal{L}}{\partial \mathbf{y}} \right)^T 
        \cdot \mathbf{h}^T.
        \]
    \end{itemize}

    \item Gradient with respect to \( \mathbf{b}_2 \):
    \[
    \frac{\partial \mathcal{L}}{\partial \mathbf{b}_2} = 
    \frac{\partial \mathcal{L}}{\partial \mathbf{y}}.
    \]
    This is because \( \mathbf{y} = W_2 \mathbf{h} + \mathbf{b}_2 \), 
    so \( \frac{\partial \mathbf{y}}{\partial \mathbf{b}_2} = I \), 
    where \( I \) is the identity matrix.
\end{itemize}

\subsection{Hidden Layer Gradients}

The gradients of the loss \( \mathcal{L} \) with respect to the 
hidden layer parameters \( W_1 \) and \( \mathbf{b}_1 \) are:

\begin{itemize}
    \item Gradient with respect to \( W_1 \):
    \[
    \frac{\partial \mathcal{L}}{\partial W_1} = 
    \left( \frac{\partial \mathcal{L}}{\partial \mathbf{h}} 
    \odot \sigma^{\prime}(\mathbf{z}) \right)^T \cdot \mathbf{x}^T.
    \]
    This is derived as follows:
    \begin{itemize}
        \item The chain rule gives:
        \[
        \frac{\partial \mathcal{L}}{\partial W_1} = 
        \frac{\partial \mathcal{L}}{\partial \mathbf{h}} 
        \cdot \frac{\partial \mathbf{h}}{\partial W_1}.
        \]
        \item Since \( \mathbf{h} = \sigma(\mathbf{z}) \) and 
        \( \mathbf{z} = W_1 \mathbf{x} + \mathbf{b}_1 \), the 
        derivative \( \frac{\partial \mathbf{h}}{\partial W_1} \) is a 
        \( 3 \times 3 \times 2 \) tensor. To compute it, we first write:
        \[
        \frac{\partial \mathcal{L}}{\partial W_1} = 
        \frac{\partial \mathcal{L}}{\partial h_1} \cdot 
        \frac{\partial h_1}{\partial W_1} + 
        \frac{\partial \mathcal{L}}{\partial h_2} \cdot 
        \frac{\partial h_2}{\partial W_1} + 
        \frac{\partial \mathcal{L}}{\partial h_3} \cdot 
        \frac{\partial h_3}{\partial W_1}.
        \]
        \item The partial derivatives \( \frac{\partial h_1}{\partial W_1} \), 
        \( \frac{\partial h_2}{\partial W_1} \), and \( \frac{\partial h_3}{\partial W_1} \) 
        are matrices of size \( 3 \times 2 \), where:
        \[
        \frac{\partial h_1}{\partial W_1} = 
        \begin{bmatrix}
        \sigma^{\prime}(z_1) x_1 & \sigma^{\prime}(z_1) x_2 \\
        0 & 0 \\
        0 & 0
        \end{bmatrix}, \quad
        \frac{\partial h_2}{\partial W_1} = 
        \begin{bmatrix}
        0 & 0 \\
        \sigma^{\prime}(z_2) x_1 & \sigma^{\prime}(z_2) x_2 \\
        0 & 0
        \end{bmatrix}, \quad
        \frac{\partial h_3}{\partial W_1} = 
        \begin{bmatrix}
        0 & 0 \\
        0 & 0 \\
        \sigma^{\prime}(z_3) x_1 & \sigma^{\prime}(z_3) x_2
        \end{bmatrix}.
        \]
        \item Substituting these into the expression for 
        \( \frac{\partial \mathcal{L}}{\partial W_1} \), we get:
        \[
        \frac{\partial \mathcal{L}}{\partial W_1} = 
        \frac{\partial \mathcal{L}}{\partial h_1} \cdot 
        \begin{bmatrix}
        \sigma^{\prime}(z_1) x_1 & \sigma^{\prime}(z_1) x_2 \\
        0 & 0 \\
        0 & 0
        \end{bmatrix} + 
        \frac{\partial \mathcal{L}}{\partial h_2} \cdot 
        \begin{bmatrix}
        0 & 0 \\
        \sigma^{\prime}(z_2) x_1 & \sigma^{\prime}(z_2) x_2 \\
        0 & 0
        \end{bmatrix} + 
        \frac{\partial \mathcal{L}}{\partial h_3} \cdot 
        \begin{bmatrix}
        0 & 0 \\
        0 & 0 \\
        \sigma^{\prime}(z_3) x_1 & \sigma^{\prime}(z_3) x_2
        \end{bmatrix}.
        \]
        \item This results in:
        \[
        \frac{\partial \mathcal{L}}{\partial W_1} = 
        \begin{bmatrix}
        \frac{\partial \mathcal{L}}{\partial h_1} \sigma^{\prime}(z_1) x_1 & 
        \frac{\partial \mathcal{L}}{\partial h_1} \sigma^{\prime}(z_1) x_2 \\
        \frac{\partial \mathcal{L}}{\partial h_2} \sigma^{\prime}(z_2) x_1 & 
        \frac{\partial \mathcal{L}}{\partial h_2} \sigma^{\prime}(z_2) x_2 \\
        \frac{\partial \mathcal{L}}{\partial h_3} \sigma^{\prime}(z_3) x_1 & 
        \frac{\partial \mathcal{L}}{\partial h_3} \sigma^{\prime}(z_3) x_2
        \end{bmatrix}.
        \]
        \item This can be written compactly as:
        \[
        \frac{\partial \mathcal{L}}{\partial W_1} = 
        \left( \frac{\partial \mathcal{L}}{\partial \mathbf{h}} 
        \odot \sigma^{\prime}(\mathbf{z}) \right)^T \cdot \mathbf{x}^T.
        \]
    \end{itemize}

    \item Gradient with respect to \( \mathbf{b}_1 \):
    \[
    \frac{\partial \mathcal{L}}{\partial \mathbf{b}_1} = 
    \frac{\partial \mathcal{L}}{\partial \mathbf{h}} \odot \sigma^{\prime}(\mathbf{z}).
    \]
    This is derived as follows:
    \begin{itemize}
        \item The chain rule gives:
        \[
        \frac{\partial \mathcal{L}}{\partial \mathbf{b}_1} = 
        \frac{\partial \mathcal{L}}{\partial \mathbf{h}} 
        \cdot \frac{\partial \mathbf{h}}{\partial \mathbf{b}_1}.
        \]
        \item Since \( \mathbf{h} = \sigma(\mathbf{z}) \) and 
        \( \mathbf{z} = W_1 \mathbf{x} + \mathbf{b}_1 \), we have:
        \[
        \frac{\partial \mathbf{h}}{\partial \mathbf{b}_1} = 
        \frac{\partial \mathbf{h}}{\partial \mathbf{z}} \cdot 
        \frac{\partial \mathbf{z}}{\partial \mathbf{b}_1}.
        \]
        \item The derivative \( \frac{\partial \mathbf{h}}{\partial \mathbf{z}} \) 
        is a diagonal matrix where each diagonal entry is \( \sigma^{\prime}(z_i) \):
        \[
        \frac{\partial \mathbf{h}}{\partial \mathbf{z}} = 
        \begin{bmatrix}
        \sigma^{\prime}(z_1) & 0 & 0 \\
        0 & \sigma^{\prime}(z_2) & 0 \\
        0 & 0 & \sigma^{\prime}(z_3)
        \end{bmatrix}.
        \]
        \item The derivative \( \frac{\partial \mathbf{z}}{\partial \mathbf{b}_1} \) 
        is the identity matrix \( I \) because \( \mathbf{z} = W_1 \mathbf{x} + \mathbf{b}_1 \):
        \[
        \frac{\partial \mathbf{z}}{\partial \mathbf{b}_1} = I.
        \]
        \item Thus, the derivative \( \frac{\partial \mathbf{h}}{\partial \mathbf{b}_1} \) 
        is:
        \[
        \frac{\partial \mathbf{h}}{\partial \mathbf{b}_1} = 
        \frac{\partial \mathbf{h}}{\partial \mathbf{z}} \cdot I = 
        \begin{bmatrix}
        \sigma^{\prime}(z_1) & 0 & 0 \\
        0 & \sigma^{\prime}(z_2) & 0 \\
        0 & 0 & \sigma^{\prime}(z_3)
        \end{bmatrix}.
        \]
        \item Finally, the gradient \( \frac{\partial \mathcal{L}}{\partial \mathbf{b}_1} \) 
        is computed as:
        \[
        \frac{\partial \mathcal{L}}{\partial \mathbf{b}_1} = 
        \frac{\partial \mathcal{L}}{\partial \mathbf{h}} \cdot 
        \frac{\partial \mathbf{h}}{\partial \mathbf{b}_1} = 
        \frac{\partial \mathcal{L}}{\partial \mathbf{h}} \odot \sigma^{\prime}(\mathbf{z}),
        \]
        where \( \odot \) denotes the element-wise (Hadamard) product.
    \end{itemize}
\end{itemize}
\end{notechapter}
